\subsection{Row Reduction and Elementary Matrices}

\begin{definition} \label{def:elementary-row-operations-and-matrices} \hlt{Elementary Row Operations and Elementary Matrices}\\
For row symbols $R_i$, the \hlt{elementary row operations} are:
\begin{enumerate}[label=(\roman*)]
\setlength{\itemsep}{0pt}
\item \hlt{Row scaling}: $R_i\leftarrow cR_i$, $c\in\mathbb{F}^{\times}$;
\item \hlt{Row swap}: $R_i\leftrightarrow R_j$, $i\neq j$;
\item \hlt{Row replacement}: $R_i\leftarrow R_i+cR_j$, $i\neq j$, $c\in\mathbb{F}$.
\end{enumerate}
An \hlt{elementary matrix} is $E=\epsilon(I_m)$, where $\epsilon$ is one elementary row operation.
\end{definition}

\begin{proposition} \label{prop:elementary-matrices-perform-row-operations} \hlt{Elementary Matrices Perform Row Operations}\\
Let $A\in\operatorname{Mat}(m,n,\mathbb{F})$.
\begin{enumerate}[label=(\roman*)]
\setlength{\itemsep}{0pt}
\item If $\epsilon$ is an elementary row operation and $E_{\epsilon}=\epsilon(I_m)$, then $E_{\epsilon}A=\epsilon(A)$.
\item $E_{\epsilon}^{-1}=E_{\epsilon^{-1}}$; hence every elementary matrix is invertible with elementary inverse.
\item If $F=E_{\epsilon}^t\in\operatorname{Mat}(n,n,\mathbb{F})$, then $AF=(\epsilon(A^t))^t$.\\
Hence right multiplication performs the corresponding column operation.
\end{enumerate}
\begin{proof}
\begin{enumerate}[label=(\roman*)]
\setlength{\itemsep}{0pt}
\item For each row index $r$, $\operatorname{row}_r(E_{\epsilon}A)=\operatorname{row}_r(E_{\epsilon})A$.\\
So the rows of $E_{\epsilon}A$ are exactly the $\epsilon$-linear combinations of the rows of $A$.
\item Since $\epsilon^{-1}$ is elementary, $E_{\epsilon^{-1}}E_{\epsilon} =E_{\epsilon^{-1}}\epsilon(I_m) =\epsilon^{-1}(\epsilon(I_m)) =I_m$, and similarly $E_{\epsilon}E_{\epsilon^{-1}}=I_m$.
\item By (i) and Proposition \ref{prop:transpose-reverses-products}, $(AF)^t=E_{\epsilon}A^t=\epsilon(A^t)$, so $AF=(\epsilon(A^t))^t$.
\end{enumerate}
\hlt{Structural Commentary}: Elementary matrices encode row and column operations by multiplication.\\
\hlt{Key Insight}: Row operations are left multiplication; column operations are right multiplication.
\end{proof}
\end{proposition}

\begin{definition} \label{def:row-equivalence-and-rref} \hlt{Row Equivalence and Reduced Row Echelon Form}\\
For $A,B\in\operatorname{Mat}(m,n,\mathbb{F})$, write $A\sim_{\mathrm{row}}B$ if $B=E_k\cdots E_1A$ for elementary $E_i\in\operatorname{Mat}(m,m,\mathbb{F})$ and some $k\geq0$.\\
A matrix $R=[r_{ij}]$ is in \hlt{reduced row echelon form} if its nonzero rows have pivot columns $p_1,\ldots,p_s$ such that:
\begin{enumerate}[label=(\roman*)]
\setlength{\itemsep}{0pt}
\item rows $1,\ldots,s$ are nonzero and rows $s+1,\ldots,m$ are zero;
\item $r_{i,p_i}=1$ and $r_{i,j}=0$ for $j<p_i$;
\item $p_1<\cdots<p_s$;
\item $r_{\ell,p_i}=0$ for all $\ell\neq i$.
\end{enumerate}
\end{definition}

\begin{theorem} \label{thm:rref-existence-uniqueness-and-elementary-factorization} \hlt{Reduced Row Echelon Form}\\
For every $A\in\operatorname{Mat}(m,n,\mathbb{F})$, there is a unique RREF matrix $\operatorname{rref}(A)$ with $A\sim_{\mathrm{row}}\operatorname{rref}(A)$.\\
If $R=\operatorname{rref}(A)$, then $R=E_k\cdots E_1A$ for elementary $E_i\in\operatorname{Mat}(m,m,\mathbb{F})$.
\begin{proof}
Gaussian elimination successively chooses a pivot column, swaps/scales to a pivot $1$, clears the pivot column, and continues on the remaining submatrix. Proposition \ref{prop:elementary-matrices-perform-row-operations} gives $R=E_k\cdots E_1A$ with $R$ in RREF.\\
For uniqueness, row operations preserve $\ker(A)=\{\mathbf{x}:A\mathbf{x}=\mathbf{0}\}$. For an RREF matrix, the pivot set and pivot-variable formulas are determined by this kernel.\\
Hence two row-equivalent RREF matrices have the same pivot columns and rows.\\
\hlt{Structural Commentary}: RREF is the canonical row-operation representative of a matrix.\\
\hlt{Key Insight}: Elimination gives existence; the solution space pins down the reduced representative.
\end{proof}
\end{theorem}

\begin{lemma} \label{lem:invertible-matrices-row-reduce-to-identity} \hlt{Invertible Matrices Row-Reduce to the Identity}\\
If $P\in\operatorname{Mat}(m,m,\mathbb{F})$ is invertible, then $\operatorname{rref}(P)=I_m$. Hence there are elementary matrices $E_1,\ldots,E_k\in\operatorname{Mat}(m,m,\mathbb{F})$ such that $E_k\cdots E_1P=I_m$.
\begin{proof}
Let $R=\operatorname{rref}(P)$. By Theorem \ref{thm:rref-existence-uniqueness-and-elementary-factorization}, $R=E_k\cdots E_1P$ for elementary matrices $E_i$. Since each $E_i$ is invertible by Proposition \ref{prop:elementary-matrices-perform-row-operations}, left multiplication by $E_k\cdots E_1$ preserves the kernel, so $\ker(R)=\ker(P)$. Because $P$ is invertible, $P\mathbf{x}=\mathbf{0}$ implies $\mathbf{x}=P^{-1}P\mathbf{x}=\mathbf{0}$; hence $\ker(R)=\{\mathbf{0}\}$.\\
Suppose the nonzero rows of $R$ have pivot columns $p_1<\cdots<p_s$. If $s<m$, choose a non-pivot column $j$. Set the non-pivot variable $x_j=1$, set all other non-pivot variables equal to $0$, and define the pivot variables recursively from $i=s$ down to $1$ so that the $i$th row equation is satisfied. This gives a nonzero $\mathbf{x}$ with $R\mathbf{x}=\mathbf{0}$, contradicting $\ker(R)=\{\mathbf{0}\}$. Thus $s=m$.\\
Since $p_1<\cdots<p_m$ are $m$ columns in an $m\times m$ matrix, $p_i=i$ for every $i$. The RREF conditions then force the $i$th pivot column to be the $i$th standard coordinate column, so $R=I_m$.
\end{proof}
\end{lemma}

\begin{corollary} \label{cor:row-equivalence-characterization} \hlt{Row Equivalence Characterization}\\
For $A,B\in\operatorname{Mat}(m,n,\mathbb{F})$, the following are equivalent:
\begin{enumerate}[label=(\roman*)]
\setlength{\itemsep}{0pt}
\item $A\sim_{\mathrm{row}}B$.
\item $\operatorname{rref}(A)=\operatorname{rref}(B)$.
\item $\exists\,P\in\operatorname{Mat}(m,m,\mathbb{F})^{\times}$ such that $B=PA$.
\end{enumerate}
\begin{proof}
\begin{enumerate}[label=(\roman*).]
\setlength{\itemsep}{0pt}
\item[$(i)$]$\Rightarrow (ii)$ Row-equivalent matrices have the same unique RREF by Theorem \ref{thm:rref-existence-uniqueness-and-elementary-factorization}.
\item[$(ii)$]$\Rightarrow (iii)$ If $R=\operatorname{rref}(A)=\operatorname{rref}(B)$, then $R=E_k\cdots E_1A=F_{\ell}\cdots F_1B$.\\
Since each $F_j$ is invertible, $B=(F_{\ell}\cdots F_1)^{-1}E_k\cdots E_1A$.
\item[$(iii)$]$\Rightarrow (i)$ If $B=PA$ and $P\in\operatorname{Mat}(m,m,\mathbb{F})^{\times}$, Lemma \ref{lem:invertible-matrices-row-reduce-to-identity} gives elementary matrices with $E_k\cdots E_1P=I_m$.\\
Hence $P=E_1^{-1}\cdots E_k^{-1}$, so $B$ is obtained from $A$ by elementary row operations.
\end{enumerate}
\hlt{Structural Commentary}: Row equivalence is invertible left multiplication.\\
\hlt{Key Insight}: Products of elementary matrices are the coordinate record of row reduction.
\end{proof}
\end{corollary}

\begin{method} \label{meth:augmented-inverse-computation-preview} \hlt{Augmented Inverse Computation Preview}\\
For $A\in\operatorname{Mat}(n,n,\mathbb{F})$:
\begin{enumerate}[label=(\roman*)]
\setlength{\itemsep}{0pt}
\item Form $[\,A\ I_n\,]$.
\item Apply Gauss--Jordan reduction to the left block.
\item If $[\,A\ I_n\,]\sim[\,I_n\ B\,]$, then $B=A^{-1}$.
\item Otherwise $A$ is not invertible, by Corollary \ref{cor:square-matrix-invertible-iff-full-rank}.
\end{enumerate}
\end{method}
